\chapter*{What you need for this course}
\markboth{Course requirements}{Course requirements}
\addcontentsline{toc}{chapter}{What you need for this course}

To follow the course and to try out the ideas of this book on a real machine, you need the following.

\begin{keyidea}[Requirements]
\begin{itemize}[nosep]
  \item \textbf{A laptop with Ubuntu 26.04.} It is used for the problems marked with \texttt{[code]}, for plotting and for programming the examples of this book.
  \item \textbf{A Booster K1 robot.} The humanoid robot from the title page and from the webster flip in the Motivation chapter. You will use it to test what you have learned on a real robot.
\end{itemize}
\end{keyidea}

\textbf{Before the first lesson.} Make sure that the laptop starts and that you can log in to your Ubuntu account. Check that the robot is charged and that you know who is responsible for it in your group.

\textbf{Working in groups.} If robots or laptops are shared, agree in your group who is using which device and when, so that everybody can try the code and the robot.
